% page 575 of the facsimile (image p-574 of the scan)
% block 170.2 x 267.9 mm ; left margin 31.8 mm
\pgc{10.160mm}{0.000mm}{
\l{}
\l{}
\l{}
\l{}
\l{\cel{53}567}
\l{}
\l{\sou{tala}.- Sama kam kozo, persono, pri qua onu jus parolis o quik}
\l{\cel{3}parolos. -}
\l{}
\l{\sou{talo}. (\sou{bot}.) La organaro vejetantala, plu o min komplexa, di la}
\l{\cel{3}planti minim supra pri organizeso : algi, fungi, e c. - DEFIS.}
\l{}
\l{\sou{talento}. Valoro, ecelo, naturala od aquirita per exercado, en}
\l{\cel{3}irga ek la generi. - DEFIRS.}
\l{}
\l{\sou{taliar}. (\sou{trans.})\cel{2}Tranchar (ulo) segun formo determinita.- EFIS.}
\l{}
\l{\sou{talio}. (\sou{kemio}).Korpo simpla, metalo min blanka kam arjento. -}
\l{\cel{3}\sou{Simbolo kemiala}\cel{1}:\cel{2}Tl.\cel{2}- DEFIRS.}
\l{}
\l{\sou{taliono}. (\sou{yuro-cienco}).Puniso qua konsistas ye aplikar a la kul-}
\l{\cel{3}pinto la ipsa trakto quan lu subisigis da sua kunhomo. - F.}
\l{}
\l{\sou{talioro}. La persono qua facas la vesti di viri e di mulieri. -EFI}
\l{}
\l{\sou{talismano}. Lapido, ringo, e c., sur qua trasesis ula signi pis-}
\l{\cel{3}terioza, e qua egardesas (da ula personi) kom apta efikar super-}
\l{\cel{3}nature pro lia relato kun la stelari sub qui li gravesis. -}
\l{\cel{3}DEFIRS.}
\l{}
\l{\sou{talko.}\cel{2}Substanco lameloza, griza verdatre, grasatra tushale.}
\l{\cel{3}\sou{Simbolo kemiala}\cel{1}: H2 Hg3 Si4 O12. - DEFIRS.}
\l{}
\l{\sou{talono}. I. (\sou{anat.}) La dopa parto di la pedo, che la homo. -}
\l{\cel{3}II. Parto di la pedo-vesto qua inkluzas la dopa parto di la}
\l{\cel{3}homo-pedo. - III. (\sou{metaf.}) (a) Extremajo infra di ula objekti;}
\l{\cel{3}(b) Marjino di la folieti detranchebla, qua duras sur parto}
\l{\cel{3}di la folieto\cel{2}fixigita ye la registro, ye la kayereto, e qua}
\l{\cel{3}uzesas kom kontrolo-moyeno. - DFIS.}
\l{}
\l{\sou{talpo.}\cel{2}(\sou{zool.)}\cel{2}Mikra mamifero karnavida, qua exkavas galerii}
\l{\cel{3}sub la sulsurfaco, e di qua la okuli esas extreme mikra. -}
\l{\cel{3}L.\cel{1}\sou{talpa}.\cel{2}FISL.}
\l{}
\l{\sou{taluso}. Tereno, kun pento tre inklinita, qua rezultas ordinare}
\l{\cel{3}de terlaboro, e qua garnisas la bordo di trancheo, di fosato,}
\l{\cel{3}e c. - DEFS.}
\l{}
\l{\sou{talvego}. (\sou{geogr.}) Lineo qua sequas la fundo di valo, di la lito}
\l{\cel{3}di fluvio, di rivero - inter-sekuro di du penti laterala.}
\l{\cel{3}- DEFI.}
\l{}
\l{\sou{tam.}\cel{2}Adverbo : sama-grada pri qualeso.}
\l{}
\l{\sou{tamarindo}.\cel{1}\sou{(bot.)}\cel{2}Frukto de la arboreto L.\cel{1}\sou{tamarindus,}\cel{1}ek la}
\l{\cel{3}familio \textquotedbl{}leguminosi\textquotedbl{}, di qua la shelo kontenas pulpo laxigiva,}
\l{\cel{3}mi-purgiva. - DEFIRS.}
\l{}
\l{tamarisko. (bot.) Arboreto di qua la folii esas tenua, e la}
\l{\cel{3}flori spikoza. - L.\cel{1}\sou{tamarix.}\cel{3}DEFIRS.}
\l{}
\l{tamburo.\cel{2}Kesto cilindra di qua la du fundi esas kovrita ye}
}
